Satellite embeddings have become a practical interface between large-scale Earth observation data and downstream geospatial analysis, yet the literature is still organized mainly around foundation models rather than the embeddings they produce. This review reframes the area from an embedding-centered perspective. We first define satellite embeddings as reusable latent representations derived from satellite or multimodal Earth observation data, and organize existing methods along three axes: model family, data and modality support, and deployment mode. We then consolidate \modelcount{} representative models into a unified dossier covering publication source, authorship, spatial scale, temporal semantics, architecture, pretraining signal, embedding dimensionality, runtime behavior, and code or artifact availability. Building on the organization of recent remote sensing reviews, we connect the model taxonomy to application requirements, data regimes, output granularity, and operational constraints. We further report a compact Wuhan benchmark that evaluates frozen embeddings on classification, fraction regression, water extraction, and dense segmentation within a shared 500\,m grid. The resulting review is intended to serve as both a technical reference and an evidence-backed evaluation framework for selecting satellite embeddings in practical Earth observation workflows.
